\documentclass[../../main.tex]{subfiles}

\begin{document}

\subsection{GRAIN} \label{subsec:grain}
  \begin{figure}[ht]
    \centering
    \includegraphics[width=.2\textwidth]{DUNE/sand/grain-front}
    \hspace{1.6cm}
    \includegraphics[width=.385\textwidth]{DUNE/sand/grain-side}
    \caption{
      Rendering frontale (sinistra) e laterale (destra) di GRAIN che ne mostra
      il recipiente interno ed esterno.
      Il recipiente interno è realizzato in acciaio inox, mentre quello esterno
      in fibra di carbonio con una struttura a nido d'ape al fine di ridurre al
      minimo la quantità di materiale all'esterno del bersaglio.
      Tra i due recipienti verrà mantenuto il vuoto per isolare
      termicamente quello interno.
    }
    \label{fig:grain}
  \end{figure}
  GRAIN (\textit{Granular Argon for Interaction of Neutrinos})
  è un bersaglio attivo ad argon liquido con una massa fiduciale di
  $\sim$1~tonnellata, che sarà posizionato nella sezione a monte del volume
  interno di SAND\@.
  Il suo obiettivo principale è vincolare gli effetti nucleari sull'argon e
  disporre di un bersaglio di argon complementare costantemente posizionato
  \textit{on-axis} rispetto al fascio di neutrini,
  per la calibrazione incrociata con gli altri rivelatori.

  Come mostrato in figura~\ref{fig:grain}, GRAIN consiste in due criostati
  cilindrici coassiali: un recipiente interno e uno esterno.
  Quello interno, conterrà l'argon liquido a una pressione di 1,5~bar, per
  isolarlo tra i due recipienti verrà mantenuto il vuoto a una pressione di
  $10^{-4} - 10^{-5}$~bar.\\
  GRAIN si basa sull'imaging della luce di scintillazione emessa dalle
  particelle cariche, per questo
  motivo, il recipiente interno sarà dotato di fotocamere
  \textit{Vacuum UltraViolet} per raccogliere la luce di scintillazione
  dell'argon.

\end{document}